\documentclass[11pt]{article}
\usepackage[margin=1in]{geometry}
\usepackage{amsmath,amssymb,amsthm}
\usepackage[colorlinks=true,linkcolor=blue,urlcolor=blue,citecolor=blue]{hyperref}
\newtheorem{theorem}{Theorem}[section]
\newtheorem{lemma}[theorem]{Lemma}
\newtheorem{proposition}[theorem]{Proposition}
\newtheorem{corollary}[theorem]{Corollary}
\newtheorem{remark}[theorem]{Remark}
\numberwithin{equation}{section}
\newcommand{\tr}{\operatorname{tr}}
\newcommand{\rank}{\operatorname{rank}}
\newcommand{\diag}{\operatorname{diag}}
\newcommand{\spec}{\operatorname{spec}}
\newcommand{\op}{\mathrm{op}}
\newcommand{\R}{\mathbb R}
\newcommand{\C}{\mathbb C}
\newcommand{\cH}{\mathcal H}
\newcommand{\cN}{\mathcal N}
\title{One rank-one perturbation and many fractional singularities:\\
sharp localization and amplitude-counting laws}
\author{Working research manuscript}
\date{November 2026}
\begin{document}
\maketitle

\begin{abstract}
For a finite sum of shifted homogeneous fractional functions, with a
common exponent $0<\alpha<1$ and an analytic remainder, we determine the
exact small-exponent Schatten power constant and the leading worst-case
approximation rank. The coefficients add over the isolated singularities,
although one positive rank-one perturbation already realizes the lower
bound. The proof glues finite local witnesses by shrinking a single
perturbation around separated scalar spectral blocks. A fixed-exponent
mass-allocation inequality yields uniformity in arbitrarily small
amplitudes, arbitrarily close distinct singularities, and a growing
number of singularities. For infinitely many isolated cusps with amplitude
counting function $A(u)\sim C u^\nu$, the optimal power constant grows as
$s^{-2-\nu}$ and the compression rank as $\log^{2+\nu}(1/\varepsilon)$,
with explicit matching constants. All input matrices in these statements
are finite and have spectra in a fixed compact interval. The update norm
may approach zero, and error is normalized by its $\alpha$ power.
The sharp single-cusp constants are inputs from a companion manuscript;
singular-support localization and separated-block perturbation methods
have classical antecedents.
\end{abstract}

\noindent\textit{Review and priority status.}
The complete mathematical note and typeset source passed a separate
independent internal audit. No external priority, human peer review,
formal verification, arithmetic runtime bound, or finite-word stability
is asserted.

\section{Definitions, statements, and antecedents}

Fix $0<\alpha<1$ and $\beta=1-\alpha$. Let $I=[a,b]$, where $a<b$.
Matrices are finite Hermitian; dimensions are unrestricted unless
specified. The scalar function may be complex valued, so its
functional-calculus output need not be Hermitian. Define the power constant
\begin{equation}\label{eq:H}
 \cH_s^I(f)=
 \sup_{\substack{\spec(A),\spec(B)\subset I\\A\ne B}}
 \frac{\tr|f(B)-f(A)|^s}{\tr|B-A|^{\alpha s}},\qquad 0<s\le1.
\end{equation}
For a fixed positive integer $r$, let $\cN_{f,r}^I(L)$ be the supremum of
\begin{equation}\label{eq:N}
 \#\{k:\sigma_k(f(B)-f(A))>e^{-L}\|B-A\|_{\op}^{\alpha}\}
\end{equation}
over the same pairs with $\rank(B-A)\le r$.
This uses a strict threshold. By the best-rank approximation theorem,
these counts describe optimal universal approximation ranks at the stated
error. They do not assert a singular-value asymptotic for a fixed pair.

Write $g_+(x)=(-ix)^\alpha$ and $g_-(x)=(ix)^\alpha$, using principal
branches and value zero at zero. For
\begin{equation}\label{eq:F}
 F=\kappa g_++\lambda g_-,\qquad \kappa,\lambda\in\C,
\end{equation}
put
\begin{equation}\label{eq:channelcount}
 m(F)=1_{\kappa\ne0}+1_{\lambda\ne0},\qquad
 c(F)=\frac{m(F)}{2\pi^2\alpha\beta}.
\end{equation}
The companion \cite{Channels} proves the exact whole-line constants
\begin{equation}\label{eq:companion}
 s^2\cH_s(F)\longrightarrow 2c(F),\qquad
 \frac{\cN_{F,r}(L)}{L^2}\longrightarrow r c(F),
\end{equation}
where the superscript $I$ is omitted for unrestricted real spectra.
Its rank-one reduction identifies $\cH_s(F)$ with the supremum of
$\tr|F(B)-F(A)|^s$ over positive rank-one updates
$B-A=vv^*$, $\|v\|=1$. Positive rank-one examples also attain the
leading count in \eqref{eq:companion}. Homogeneity identifies the
intrinsic error scale in \eqref{eq:N} with unit-update normalization.
The same companion gives
\begin{equation}\label{eq:uniformsingle}
 \sup_{(\kappa,\lambda)\ne(0,0)}
 \left|\frac{s^2\cH_s(F)}{|\kappa|^s+|\lambda|^s}
                  -\frac1{\pi^2\alpha\beta}\right|\longrightarrow0
\end{equation}
and the operator estimate
$\|F(B)-F(A)\|_{\op}\le C_F\|B-A\|_{\op}^{\alpha}$.
These are explicit inputs, rather than new single-cusp claims here.

Let $c_1,\ldots,c_M$ be distinct points in the interior of $I$. Consider
\begin{equation}\label{eq:finitef}
 f(x)=h(x)+\sum_{j=1}^M F_j(x-c_j),
\end{equation}
where $h$ is holomorphic on a complex neighborhood of $I$ and every
$F_j$ has the form \eqref{eq:F}, with the same $\alpha$.
Zero terms are omitted. Put $Q=\sum_j m(F_j)$.

\begin{theorem}[Finite singularities]\label{thm:finite}
For \eqref{eq:finitef},
\begin{equation}\label{eq:finiteH}
 \lim_{s\downarrow0}s^2\cH_s^I(f)
                  =\frac{Q}{\pi^2\alpha\beta}.
\end{equation}
For every fixed positive integer $r$,
\begin{equation}\label{eq:finiteN}
 \lim_{L\to\infty}\frac{\cN_{f,r}^I(L)}{L^2}
                  =\frac{rQ}{2\pi^2\alpha\beta}.
\end{equation}
Positive rank-one increments attain the leading lower constants for
$r=1$; their $r$-fold direct sums attain \eqref{eq:finiteN}.
The locations and nonzero fixed amplitudes do not enter these constants.
\end{theorem}

Every nonzero real homogeneous fractional function
$u x_+^\alpha+v(-x)_+^\alpha$ has $m(F)=2$.
Consequently a nonzero real coefficient at each of $M$ shifted fractional
ReLU singularities gives $Q=2M$. Its moment coefficient is
$2M/(\pi^2\alpha\beta)$ and its count coefficient is
$rM/(\pi^2\alpha\beta)$, even after adding an analytic $h$.
The companion PSD power coefficient is $1/(2\pi^2\alpha\beta)$
\cite{Sharp}, so the power-constant multiplier is $4M$.
Ordinary Schatten norm constants are the $1/s$ powers of these power
constants; the comparison is not an ordinary norm ratio of $4M$.

\begin{theorem}[Fixed-exponent mass bounds]\label{thm:mass}
For every $0<s\le1$ and positive masses $p_j$ summing to one,
\begin{equation}\label{eq:mass}
 \sum_j p_j^{\alpha s}\cH_s(F_j)
 \le\cH_s^I(f)
 \le\sum_j\cH_s(F_j)+\cH_s^I(h).
\end{equation}
Moreover $\cH_s^I(h)=O_{h,I}(s^{-1})$. Maximizing the lower expression
over the masses gives
\begin{equation}\label{eq:optimized}
 \cH_s^I(f)\ge
 \left[\sum_j\cH_s(F_j)^{1/(1-\alpha s)}\right]^{1-\alpha s}.
\end{equation}
This is a lower bound, rather than an assertion of exact fixed-$s$
additivity.
\end{theorem}

The update norm in \eqref{eq:H}--\eqref{eq:N} may be arbitrarily small.
This matters even on compact $I$. The separate fixed-endpoint-to-update
ratio theorem \cite{Bounded} gives the PSD count coefficient
$r/(4\pi^2\alpha)$ for any fixed ratio bound greater than one; the
unrestricted coefficient is $r/(4\pi^2\alpha\beta)$.
Our shrinking update construction recovers the unrestricted coefficient
near each interior singularity. It does not assert the same law with a
positive lower bound on the update norm.

\paragraph{Antecedents and scope.}
Singular-support localization is an established principle.
Pushnitski--Yafaev \cite[Theorem 1.1]{PY} explicitly add the contributions
of disjoint symbol singularities to power-scale Hankel singular-value
asymptotics. Arazy--Zelenko \cite[Introduction and Section 2]{AZ}
connect rank-one functional differences with Hankel and Riesz channels.
Functions with finitely many nonsmooth points already occur in
Sobolev's functional-calculus estimates \cite{Sobolev}.
Operator H\"older functions for quasi-Banach Schatten classes also have
a substantial theory \cite{HSZ}. We claim neither a new Weyl method nor
first existence of broad classes of such functions. The construction
below uses a standard first-order separated-block perturbation expansion
to glue sharp nonlinear finite-matrix witnesses while retaining one
rank-one increment. The exact extremal constants and later amplitude
laws require the companion's sharp single-cusp input. No exhaustive
priority conclusion is drawn from this bounded literature comparison.

\section{Analytic functions contribute a lower-order term}

\begin{lemma}\label{lem:analytic}
If $h$ is holomorphic on a complex neighborhood of $I$, there are
$C,c>0$ and a fixed enlarged real interval $I_\delta$ such that for
every integer $q\ge0$ a scalar polynomial $p_q$ of degree at most $q$
satisfies
\begin{equation}\label{eq:analyticerror}
 \|[h(B)-h(A)]-[p_q(B)-p_q(A)]\|_{\op}
          \le C e^{-cq}\|B-A\|_{\op}
\end{equation}
for spectra in $I_\delta$. Hence rank-one increments of norm $t$ obey
\begin{equation}\label{eq:analytictrace}
 \tr|h(B)-h(A)|^s\le\frac{(Ct)^s}{1-e^{-cs}}.
\end{equation}
For general increments with endpoints in $I$,
\begin{equation}\label{eq:analyticH}
 \cH_s^I(h)=O_{h,I}(s^{-1}).
\end{equation}
For fixed rank $r$, an analytic correction at error
$e^{-L}\|B-A\|_{\op}^{\alpha}$ requires only $O_{h,I}(r(L+1))$ rank.
\end{lemma}
\begin{proof}
Choose three nested Bernstein ellipses after affine scaling of
$I_\delta$, all inside the holomorphy domain. Let the middle ellipse
$\Gamma$ be a positive distance from $I_\delta$. Chebyshev truncation
gives $\sup_\Gamma|h-p_q|\le C e^{-cq}$.
Cauchy's formula and the resolvent identity give
\[
 [h-p_q](B)-[h-p_q](A)
 =\frac1{2\pi i}\int_\Gamma [h(z)-p_q(z)]
       (zI-B)^{-1}(B-A)(zI-A)^{-1}\,dz.
\]
Both resolvents have uniform bounds on this fixed contour, proving
\eqref{eq:analyticerror}. The derivative-scale factor $\|B-A\|_{\op}$
is essential; an absolute scalar approximation alone would not suffice.

For a polynomial $p_q$,
\begin{equation}\label{eq:polynomialrank}
 \rank(p_q(B)-p_q(A))\le q\rank(B-A).
\end{equation}
Indeed
\[
 B^k-A^k=\sum_{\ell=0}^{k-1}B^\ell(B-A)A^{k-1-\ell},
\]
and the complete column range lies in the sum of the column spaces of
$B^\ell(B-A)$ for $0\le\ell<q$.
For rank one, \eqref{eq:analyticerror} therefore bounds the $(q+1)$-st
singular value by $Ct e^{-cq}$. Summing this geometric sequence to
the $s$ power proves \eqref{eq:analytictrace}.

To prove \eqref{eq:analyticH}, put $D=b-a$, $E=B-A$, and choose a fixed
integer $K$ with $D/K<\delta$. The straight-line points $A+jE/K$ have
spectra in $I$. In each of the $K$ segments telescope the signed
spectral rank-one decomposition of $E/K$. Every partial spectral sum
has norm at most $\|E\|_{\op}/K\le D/K$, so the intermediate endpoints
have spectra in $I_\delta$. The power triangle and
\eqref{eq:analytictrace} give
\[
 \tr|h(B)-h(A)|^s
 \le\frac{C^s K^{1-s}}{1-e^{-cs}}\tr|E|^s
 \le\frac{C^s K^{1-s}D^{(1-\alpha)s}}{1-e^{-cs}}
                                            \tr|E|^{\alpha s}.
\]
This is $O(s^{-1})$ uniformly over the inputs.
Finally $\|E\|_{\op}\le D$, so \eqref{eq:analyticerror} reaches the
intrinsic error $e^{-L}\|E\|_{\op}^{\alpha}$ at $q=O_{h,I}(L+1)$;
\eqref{eq:polynomialrank} supplies the rank statement.
\end{proof}

\section{A finite-dimensional blow-up preserving rank one}

The next lemma separates the lower proof from the approximation upper.
It also applies later when only a finite selection of infinitely many
singularities is used.

\begin{lemma}[Local block limit]\label{lem:blowup}
Suppose a function $f$ is operator $\alpha$-H\"older on a neighborhood
of a compact interval $I$ and, at distinct interior points $c_j$,
has expansions
\begin{equation}\label{eq:local}
 f(c_j+x)=f(c_j)+F_j(x)+O_j(|x|),
\end{equation}
where each $F_j$ is homogeneous of degree $\alpha$.
Let $C=\diag(c_jI_{n_j})$, and fix finite Hermitian matrices
$T_+,T_-$ with this block decomposition. Then
\begin{equation}\label{eq:blowup}
 t^{-\alpha}[f(C+tT_+)-f(C+tT_-)]
 \longrightarrow
 \bigoplus_j [F_j((T_+)_{jj})-F_j((T_-)_{jj})]
\end{equation}
in operator norm as $t\downarrow0$. Both endpoint spectra belong to
$I$ for sufficiently small $t$.
\end{lemma}
\begin{proof}
For one fixed Hermitian $T$, choose an anti-Hermitian $S$ by
\[
 S_{ij}=-\frac{T_{ij}}{c_i-c_j}\quad(i\ne j),\qquad S_{ii}=0.
\]
For $U_t=e^{tS}$, direct expansion gives
\[
 U_t^*(C+tT)U_t=C+t\diag(T_{jj})+O(t^2),
\]
since $T+[C,S]$ is block diagonal. All constants here may depend on
the fixed finite blocks and the distinct centers.
The operator H\"older assumption changes this remainder by
$O(t^{2\alpha})$ in functional calculus. Undoing $U_t=I+O(t)$
changes the bounded $f$-output by $O(t)$. Thus
\[
 f(C+tT)=\bigoplus_j f(c_jI+tT_{jj})+O(t^{2\alpha}+t).
\]
Finite spectral calculus in each fixed block and \eqref{eq:local}
give
\[
 f(c_jI+tT_{jj})=f(c_j)I+t^\alpha F_j(T_{jj})+O(t).
\]
The constant blocks cancel between $T_+$ and $T_-$. Both
$2\alpha>\alpha$ and $1>\alpha$, proving \eqref{eq:blowup}.
The centers are interior and all first-order matrices are finite, so
the endpoint spectra stay in $I$ once $t$ is small enough.
\end{proof}

For \eqref{eq:finitef}, the hypotheses of Lemma~\ref{lem:blowup} hold:
other singularities are analytic near $c_j$, and the companion bounds
plus Lemma~\ref{lem:analytic} give the operator H\"older estimate on
fixed compact neighborhoods.
Take arbitrary finite positive rank-one unit-update pairs
\[
 B_j-A_j=v_jv_j^*,\qquad \|v_j\|=1,
\]
and positive masses $p_j$ summing to one. Set
\begin{equation}\label{eq:gluing}
 H_j=\frac{A_j+B_j}{2},\qquad
 v=\bigoplus_j\sqrt{p_j}\,v_j,\qquad
 T_\pm=\diag(p_j H_j)\pm\frac{vv^*}{2}.
\end{equation}
Then $\|v\|=1$. The actual global endpoints
\begin{equation}\label{eq:globalpair}
 A_t=C+tT_-,\qquad B_t=C+tT_+
\end{equation}
have $B_t-A_t=t vv^*$, one positive rank-one increment with norm exactly
$t$. Their diagonal first-order blocks are $p_jA_j$ and $p_jB_j$.
Lemma~\ref{lem:blowup} gives
\begin{equation}\label{eq:gluedlimit}
 t^{-\alpha}[f(B_t)-f(A_t)]
 \longrightarrow
 \bigoplus_j p_j^\alpha[F_j(B_j)-F_j(A_j)].
\end{equation}
This is not a direct sum of $M$ independent perturbations. Its global
increment has rank one throughout; the off-block effect vanishes
after the proved change of basis and blow-up.

\section{Finite-cusp constants and uniformity}

\begin{proof}[Proof of Theorem~\ref{thm:mass}]
The upper in \eqref{eq:mass} follows by the Schatten power triangle,
and the analytic estimate is Lemma~\ref{lem:analytic}.
For the lower choose independent finite rank-one near-maximizers
for each $\cH_s(F_j)$, using the companion reduction.
In \eqref{eq:gluedlimit} the normalized output's power trace converges,
since the matrices are finite at this point. The denominator in
\eqref{eq:H} is exactly $t^{\alpha s}$.
Let the near-maximizer errors vanish to obtain the lower in
\eqref{eq:mass}. No supremum is interchanged with an integral.
The concave maximization of $\sum_j H_j p_j^{\alpha s}$,
with $H_j=\cH_s(F_j)$ and $\sum_jp_j=1$, yields
\eqref{eq:optimized}; here $\alpha s<1$.
\end{proof}

\begin{proof}[Proof of Theorem~\ref{thm:finite}]
Fix positive masses. Multiply \eqref{eq:mass} by $s^2$ and use
\eqref{eq:companion}, $p_j^{\alpha s}\to1$, and
$s^2\cH_s^I(h)\to0$. This gives \eqref{eq:finiteH}.

For the count upper approximate every homogeneous correction at error
$e^{-L}\|B-A\|_{\op}^{\alpha}/(M+1)$, and use
Lemma~\ref{lem:analytic} for $h$.
The ranks add, while these fixed allocations only shift $L$ by constants
in \eqref{eq:companion}. This proves the upper in \eqref{eq:finiteN}.

For the $r=1$ lower fix masses and choose $C_0$ with
$p_j^\alpha e^{C_0}>2$ for all $j$.
For a given large $L$, choose independent unit positive rank-one
witnesses for $F_j$ with at least
$(c(F_j)-o(1))(L-C_0)^2$ singular values above $e^{-(L-C_0)}$.
The blocks on the right of \eqref{eq:gluedlimit} have these singular
values above $2e^{-L}$. Choose $t>0$ small enough for its operator-norm
error to be less than $e^{-L}$, preserving the strict count at $e^{-L}$.
Summing gives $\sum_jc(F_j)$ as a lower coefficient.
Take $r$ direct copies of this rank-one global construction. The total
increment has rank $r$ and still norm $t$, proving the general lower.
\end{proof}

\begin{corollary}[Many varying singularities]\label{cor:uniform}
Set $h=0$. Let $M=M_s$ be positive and finite, let all
$c_{j,s}$ be distinct interior points of $I$, and allow arbitrary
nonzero coefficient pairs $(\kappa_{j,s},\lambda_{j,s})$.
If $s\log M_s\to0$, then for
$f_s(x)=\sum_{j=1}^{M_s}F_{j,s}(x-c_{j,s})$,
\begin{equation}\label{eq:manyuniform}
 \frac{s^2\cH_s^I(f_s)}
 {\sum_j(|\kappa_{j,s}|^s+|\lambda_{j,s}|^s)}
             \longrightarrow\frac1{\pi^2\alpha\beta},
\end{equation}
uniformly in these coefficients and distinct locations.
\end{corollary}
\begin{proof}
Equal masses in \eqref{eq:mass} give
\[
 M_s^{-\alpha s}\sum_j\cH_s(F_{j,s})
 \le\cH_s^I(f_s)\le\sum_j\cH_s(F_{j,s}).
\]
Apply \eqref{eq:uniformsingle} to the sum. The loss
$M_s^{-\alpha s}$ tends to one under the stated assumption.
\end{proof}

No minimum separation, lower amplitude, or lower update scale is needed
in Corollary~\ref{cor:uniform}. The finite local witnesses are chosen
first and the global $t$ can then be arbitrarily small.
Exact coalescence can cancel a singularity and change the constant;
distinct but extremely close centers still obey the stated law.
An analytic remainder may be included with the quantified
$O(s^{-1})$ additive bound. A relative uniform formulation would also
need its size controlled against the denominator in
\eqref{eq:manyuniform}.

\section{Infinitely many cusps and amplitude counting}

Choose distinct $c_j$ in the interior of $I$, with no accumulation
point inside $I$. For example, $I=[0,1]$ and $c_j=2^{-j}$ suffice.
Let $b_j\ge1$ tend to infinity and assume
\begin{equation}\label{eq:amplitudecount}
 A(u)=\#\{j:b_j\le u\}\sim C u^\nu,
 \qquad u\to\infty,\qquad C>0,\ \nu>0.
\end{equation}
Define
\begin{equation}\label{eq:infinitef}
 f(x)=\sum_{j\ge1}e^{-b_j}|x-c_j|^\alpha.
\end{equation}
The coefficients are summable, since $A(u)\le C'(1+u^\nu)$.
The series converges locally uniformly on $\R$ and is globally
$\alpha$-H\"older both scalarly and in operator norm, with constants
bounded by the corresponding $|x|^\alpha$ constants times
$\sum_je^{-b_j}$.

\begin{theorem}[Amplitude-counting law]\label{thm:infinite}
For \eqref{eq:infinitef}, $\cH_s^I(f)<\infty$ for every $0<s\le1$ and
\begin{equation}\label{eq:infiniteH}
 \cH_s^I(f)\sim
 \frac{2C\Gamma(\nu+1)}{\pi^2\alpha\beta}\,s^{-\nu-2},
 \qquad s\downarrow0.
\end{equation}
For every fixed positive integer $r$,
\begin{equation}\label{eq:infiniteN}
 \cN_{f,r}^I(L)\sim
 \frac{2rC}{\pi^2\alpha\beta(\nu+1)(\nu+2)}L^{\nu+2},
 \qquad L\to\infty.
\end{equation}
Positive rank-one increments attain the leading lower bounds when
$r=1$. Inputs in these suprema are finite, even though $f$ has
infinitely many singularities.
\end{theorem}

Every $c_j$ is isolated. A disk around it avoiding all other centers
admits holomorphic branches of $|x-c_k|^\alpha$ for all $k\ne j$,
bounded uniformly in $k$. Their weighted series converges normally.
Thus the local expansion at $c_j$ is
\[
 f(c_j+x)=f(c_j)+e^{-b_j}|x|^\alpha+O_j(|x|).
\]
Together with the global operator H\"older bound, these are exactly the
hypotheses of Lemma~\ref{lem:blowup} for any finite selected set of
centers. The unselected cusps need not form a globally analytic remainder.

\subsection{Power constants}

Put $F_0(x)=|x|^\alpha$ and
\[
 c_0=\frac1{\pi^2\alpha\beta},\qquad
 H_0(s)=\cH_s(F_0)\sim\frac{2c_0}{s^2},\qquad
 D(s)=\sum_j e^{-s b_j}.
\]
The infinite power triangle gives
\begin{equation}\label{eq:infiniteupper}
 \cH_s^I(f)\le H_0(s)D(s).
\end{equation}
It is justified by finite partial sums followed by their norm limit
on each fixed finite matrix pair; the power trace is continuous there.
On every finite selected set, the blow-up and optimized masses give
the lower bound of Theorem~\ref{thm:mass}. Increasing the selected set
yields
\begin{equation}\label{eq:infinitelower}
 \cH_s^I(f)\ge
 H_0(s)D\left(\frac{s}{1-\alpha s}\right)^{1-\alpha s}.
\end{equation}
All sums are finite for positive $s$. Stieltjes integration gives
\begin{equation}\label{eq:D}
 D(s)=s\int_0^\infty e^{-su}A(u)\,du
                  \sim C\Gamma(\nu+1)s^{-\nu}.
\end{equation}
After $v=su$, dominated convergence applies using
$A(u)\le C'(1+u^\nu)$. If $q_s=s/(1-\alpha s)$, then $q_s/s\to1$ and
\begin{equation}\label{eq:Dsqueeze}
 \frac{D(q_s)^{1-\alpha s}}{D(s)}\longrightarrow1,
\end{equation}
because $\log D(s)=\nu\log(1/s)+O(1)$ and
$s\log D(s)\to0$. Equations
\eqref{eq:infiniteupper}--\eqref{eq:Dsqueeze} prove
\eqref{eq:infiniteH}.

\subsection{Compression upper}

Let $E=B-A$ and $t=\|E\|_{\op}$.
The polynomial bound on $A(u)$ gives
\begin{equation}\label{eq:tail}
 \sum_{b_j>T}e^{-b_j}\le C''e^{-T}(1+T^{\nu+1}).
\end{equation}
Choose
$T=L+(\nu+4)\log(L+2)$, enlarged by a fixed constant as necessary,
and retain the $J=A(T)=O(L^\nu)$ terms with $b_j\le T$.
The discarded correction has norm at most $e^{-L}t^\alpha/2$ by
\eqref{eq:tail} and the $F_0$ operator H\"older estimate.

Allocate error $e^{-L}t^\alpha/(2J)$ to each retained weighted term,
returning zero whenever this exceeds its uniform norm bound.
For every $\eta>0$, the companion count upper shows that the total
approximating rank is at most
\begin{equation}\label{eq:sumrank}
 (rc_0+\eta)\sum_j[L+O(\log L)-b_j]_+^2+O_\eta(J),
\end{equation}
uniformly over the inputs. Bounded or slightly negative effective
depths cost only a fixed finite-threshold count each and are absorbed
into $O_\eta(J)$. This also covers a companion norm bound greater than
one. The constants may depend on the fixed function parameters and
$\eta$, not the input dimension or spectra.

For $V\to\infty$, Stieltjes integration and
\eqref{eq:amplitudecount} yield
\begin{equation}\label{eq:phasevolume}
 \sum_j(V-b_j)_+^2
 =2\int_0^V(V-u)A(u)\,du
 \sim\frac{2C}{(\nu+1)(\nu+2)}V^{\nu+2}.
\end{equation}
Changing $V=L$ to $L+O(\log L)$ changes the leading term by
$o(L^{\nu+2})$. Equations \eqref{eq:sumrank}--\eqref{eq:phasevolume},
then $\eta\downarrow0$, prove the upper in \eqref{eq:infiniteN}.

\subsection{Compression lower}

For $r=1$ fix $0<\delta<1$ and retain the
$J=A((1-\delta)L)$ centers with $b_j\le(1-\delta)L$.
Assign equal masses $p_j=1/J$. Their effective depths
\begin{equation}\label{eq:effectivedepth}
 u_j=L-b_j-\alpha\log J-\log4
\end{equation}
are at least $\delta L-O(\log L)$ and tend to infinity uniformly.
Choose independent finite unit positive rank-one witnesses for $F_0$
at these depths. For all large $L$ their counts are at least
$(c_0-\eta)u_j^2$.
After multiplying by $e^{-b_j}p_j^\alpha$, the selected singular
values exceed $4e^{-L}$.

Apply Lemma~\ref{lem:blowup} to this finite selection and choose a
strictly positive $t$ sufficiently small that the normalized actual
output differs from its direct-sum limit by less than $e^{-L}$.
It follows that the actual count at $e^{-L}$ is at least
\begin{equation}\label{eq:lowercount}
 (c_0-\eta)\sum_{b_j\le(1-\delta)L}
              [L-b_j-\alpha\log J-\log4]^2.
\end{equation}
The logarithmic shifts are lower order. Integration by parts on
$[0,(1-\delta)L]$, retaining its endpoint term, gives
\begin{equation}\label{eq:truncatedvolume}
 L^{-\nu-2}\sum_{b_j\le(1-\delta)L}(L-b_j)^2
 \longrightarrow
 C\nu\int_0^{1-\delta}(1-v)^2v^{\nu-1}\,dv.
\end{equation}
For example the left side before its limit equals
\[
 \delta^2 L^{-\nu}A((1-\delta)L)
 +2\int_0^{1-\delta}(1-v)L^{-\nu}A(Lv)\,dv,
\]
so dominated convergence proves \eqref{eq:truncatedvolume} even for
$0<\nu<1$. Letting $\delta\downarrow0$ and $\eta\downarrow0$ gives
$2C c_0/[(\nu+1)(\nu+2)]$.
This is the $r=1$ lower in \eqref{eq:infiniteN}; $r$ direct copies
give the factor $r$. This proof never infers a supremal count from a
moment asymptotic by an unjustified Tauberian interchange.

\begin{corollary}[Any power beyond two]\label{cor:gamma}
For $\gamma>0$, let
$f_\gamma(x)=\sum_{j\ge1}e^{-j^\gamma}|x-c_j|^\alpha$ with the centers
as above. Then
\begin{align}
 \cH_s^I(f_\gamma)&\sim
 \frac{2\Gamma(1+1/\gamma)}{\pi^2\alpha\beta}\,
                         s^{-2-1/\gamma},\label{eq:gammaH}\\
 \cN_{f_\gamma,r}^I(L)&\sim
 \frac{r}{\pi^2\alpha\beta}
 \frac{2\gamma^2}{(\gamma+1)(2\gamma+1)}L^{2+1/\gamma}.
 \label{eq:gammaN}
\end{align}
\end{corollary}
\begin{proof}
Here $A(u)=\lfloor u^{1/\gamma}\rfloor$ asymptotically, so
$C=1$, $\nu=1/\gamma$ in Theorem~\ref{thm:infinite}.
\end{proof}

Any prescribed exponent $p>2$ results from $\gamma=1/(p-2)$.
Nevertheless $f_\gamma$ is globally $\alpha$-H\"older and has finite
$\cH_s^I$ for every positive $s\le1$.
For $\gamma=1$, the rank-one count coefficient is $c_0/3$ at $L^3$,
and the moment coefficient is $2c_0$ at $s^{-3}$.
The factor $\Gamma(4)=6$ in the layer-cake integral checks these
constants: $(c_0/3)6=2c_0$.

\section{Limits of the result}

The finite upper assumes the actual decomposition
\eqref{eq:finitef}, with analytic remainder.
A scalar $o(|x-c|^\alpha)$ remainder alone does not justify that upper.
Analytic amplitudes multiplying fractional terms, logarithmic local
multipliers, and higher-order branch corrections would need a further
sharp rational or localization argument.
Different local exponents are also not included: at update scale $t$,
normalization by $t^\alpha$ introduces factors $t^{\alpha_j-\alpha}$.

The counts are existence results for optimal low-rank approximation.
They do not supply an arithmetic algorithm with controlled pole
preparation, a lower bound on the update scale, a finite-word theorem,
or a general $N^2$ matrix algorithm. Close centers may require an
extremely small $t$ in the lower construction. The fixed positive
endpoint-to-update ratio theorem \cite{Bounded} concerns a different
class and remains distinct.

\begin{thebibliography}{9}
\bibitem{Channels}
Working companion manuscript,
\emph{Analytic channels and exact constants for homogeneous fractional
matrix maps}, November 2026.
\href{fractional-channel-classification.pdf}{Local manuscript}.

\bibitem{Sharp}
Working companion manuscript,
\emph{Exact small-exponent constants for fractional matrix maps},
November 2026.
\href{sharp-fractional-holder.pdf}{Local manuscript}.

\bibitem{Bounded}
Working research note,
\emph{Sharp fractional compression with a bounded endpoint-to-update
ratio}, November 2026.
\href{bounded-spectrum-fractional-counts.txt}{Local note}.

\bibitem{PY}
A. Pushnitski and D. Yafaev,
\emph{Localization principle for compact Hankel operators},
J. Funct. Anal. 270(9) (2016),3591--3621.
\url{https://arxiv.org/abs/1508.04279}.

\bibitem{AZ}
J. Arazy and L. Zelenko,
\emph{Directional operator differentiability of non-smooth functions},
J. Operator Theory 55(1) (2006),49--90.
\url{https://jot.theta.ro/jot/archive/2006-055-001/2006-055-001-003.pdf}.

\bibitem{HSZ}
J. Huang, F. Sukochev, and D. Zanin,
\emph{Operator $\theta$-H\"older functions with respect to
$\|\cdot\|_p$, $0<p\le\infty$},
J. London Math. Soc. (2022).
\url{https://arxiv.org/abs/2110.09708}.

\bibitem{Sobolev}
A. V. Sobolev,
\emph{Functions of self-adjoint operators in ideals of compact operators},
J. London Math. Soc. 95 (2017),157--176.
\url{https://arxiv.org/abs/1504.07261}.
\end{thebibliography}
\end{document}
